Els raigs X es generen mitjançant un tub on s'ha fet el buit i que conté dues plaques de tungstè. La placa negativa (càtode) s'escalfa i emet electrons, els quals són accelerats cap a la placa positiva (ànode). En xocar contra aquesta altra placa, els electrons s'aturen bruscament i es generen els raigs X.

\begin{apartat}{1,25}
La diferència de potencial entre dues plaques d'un tub de raigs X és de 70,0~kV i la separació entre elles és de 2,00~cm. Calculeu el mòdul del camp elèctric (suposeu que és constant i uniforme entre les plaques) i, tot seguit, el mòdul de la força elèctrica que experimenta un electró dins del tub. A la figura adjunta, representeu el camp elèctric entre plaques i la força elèctrica que actua sobre l'electró.
\end{apartat}

\begin{apartat}{1,25}
Determineu l'energia cinètica i la velocitat d'un electró quan arriba a l'ànode, suposant que parteix del repòs des del càtode. A continuació, sabent que en el xoc amb l'ànode l'electró perd part de la seva energia, que es transforma en un fotó de raigs X amb una longitud d'ona d'1,00~nm, calculeu l'energia d'aquest fotó i la nova velocitat de l'electró després de l'emissió.
\end{apartat}

\figura[0.42]{fig1.png}

\dades{%
  $|e| = 1,602\times 10^{-19}\ \mathrm{C}$.\\
  $m_\mathrm{e} = 9,11\times 10^{-31}\ \mathrm{kg}$.\\
  $c = 3,00\times 10^{8}\ \mathrm{m\,s^{-1}}$.\\
  $h = 6,63\times 10^{-34}\ \mathrm{J\,s}$.}
\nota{Negligiu els efectes del camp magnètic terrestre i els efectes relativistes.}
